\documentclass[../../main.tex]{subfiles}

\begin{document}

\section{一般可测函数的积分}\label{section:ch18-general-integral}

令 $(X,\mathcal{M},\mu)$ 为测度空间. 对 $X$ 上的可测函数 $f$, 定义它的正部 $f^{+}=\max\{f,0\}$ 和负部 $f^{-}=\max\{-f,0\}$. 则 $f^{+}$ 和 $f^{-}$ 是非负可测函数且
\begin{align*}
f=f^{+}-f^{-},\quad|f|=f^{+}+f^{-}.
\end{align*}
$f$ 称为关于 $\mu$ 可积的, 若 $\int_{X}f^{+}\,d\mu<\infty$ 且 $\int_{X}f^{-}\,d\mu<\infty$. 在这种情形下, $f$ 在 $X$ 上关于 $\mu$ 的积分定义为
\begin{align}
\int_{X}f\,d\mu=\int_{X}f^{+}\,d\mu-\int_{X}f^{-}\,d\mu.\label{align:ch18-integral-def}
\end{align}
对可测集 $E$, 定义 $\int_{E}f\,d\mu=\int_{X}f\cdot\chi_{E}\,d\mu$.

观察到 $f$ 可积当且仅当 $\int_{X}|f|\,d\mu<\infty$, 因为 $|f|=f^{+}+f^{-}$.

\begin{theorem}[积分的线性性]\label{theorem:ch18-linearity}
令 $(X,\mathcal{M},\mu)$ 为测度空间, 而 $f$ 和 $g$ 是 $X$ 上的可积函数. 则对任何实数 $\alpha$ 和 $\beta$, $\alpha f+\beta g$ 是可积的且
\begin{align*}
\int_{X}(\alpha f+\beta g)\,d\mu=\alpha\int_{X}f\,d\mu+\beta\int_{X}g\,d\mu.
\end{align*}
\end{theorem}

\begin{proof}
首先考虑 $\alpha=\beta=1$ 的情形. 令 $h=f+g$. 则 $h^{+}-h^{-}=f^{+}-f^{-}+g^{+}-g^{-}$, 因此 $h^{+}+f^{-}+g^{-}=h^{-}+f^{+}+g^{+}$. 三个非负函数的和的积分等于积分的和. 因此
\begin{align*}
\int_{X}h^{+}\,d\mu+\int_{X}f^{-}\,d\mu+\int_{X}g^{-}\,d\mu=\int_{X}h^{-}\,d\mu+\int_{X}f^{+}\,d\mu+\int_{X}g^{+}\,d\mu.
\end{align*}
所有项都是有限的, 因为 $f$ 和 $g$ 是可积的. 重新排列得到
\begin{align*}
\int_{X}h^{+}\,d\mu-\int_{X}h^{-}\,d\mu=\left[\int_{X}f^{+}\,d\mu-\int_{X}f^{-}\,d\mu\right]+\left[\int_{X}g^{+}\,d\mu-\int_{X}g^{-}\,d\mu\right].
\end{align*}
因此 $\int_{X}(f+g)\,d\mu=\int_{X}f\,d\mu+\int_{X}g\,d\mu$.

对于 $\alpha=-1$, $(-f)^{+}=f^{-}$ 且 $(-f)^{-}=f^{+}$, 因此 $\int_{X}(-f)\,d\mu=-\int_{X}f\,d\mu$. 对 $\alpha>0$, $(\alpha f)^{+}=\alpha f^{+}$ 且 $(\alpha f)^{-}=\alpha f^{-}$, 因此 $\int_{X}(\alpha f)\,d\mu=\alpha\int_{X}f\,d\mu$. 这些特殊情形的组合给出一般的线性性.
\end{proof}

\begin{proposition}[积分的比较判别法]\label{proposition:ch18-comparison}
令 $(X,\mathcal{M},\mu)$ 为测度空间. 若 $f$ 和 $g$ 是可积的且在 $X$ 上 a.e. $f\leqslant g$, 则
\begin{align*}
\int_{X}f\,d\mu\leqslant\int_{X}g\,d\mu.
\end{align*}
此外,
\begin{align*}
\left|\int_{X}f\,d\mu\right|\leqslant\int_{X}|f|\,d\mu.
\end{align*}
\end{proposition}

\begin{proof}
由于 $g-f\geqslant 0$ a.e., $\int_{X}(g-f)\,d\mu\geqslant 0$. 根据线性性, $\int_{X}g\,d\mu-\int_{X}f\,d\mu\geqslant 0$. 第二个不等式从 $-|f|\leqslant f\leqslant|f|$ 与线性性得到.
\end{proof}

以下定理是关于积分的最重要的收敛定理.

\begin{theorem}[Lebesgue 控制收敛定理]\label{theorem:ch18-dct}
令 $(X,\mathcal{M},\mu)$ 为测度空间, 而 $\{f_{n}\}$ 是 $X$ 上的可测函数序列, 满足在 $X$ 上 a.e. 逐点地 $\{f_{n}\}\to f$. 假设存在 $X$ 上的可积非负函数 $g$, 使得对所有 $n$, 在 $X$ 上 a.e. $|f_{n}|\leqslant g$. 则 $f$ 是可积的且
\begin{align}
\lim_{n\to\infty}\int_{X}f_{n}\,d\mu=\int_{X}f\,d\mu.\label{align:ch18-dct}
\end{align}
\end{theorem}

\begin{proof}
由于 $|f_{n}|\leqslant g$ a.e., 取极限得到 $|f|\leqslant g$ a.e., 因此 $f$ 是可积的. 定义 $h_{n}=g-f_{n}$ 和 $k_{n}=g+f_{n}$. 这两个函数都是非负的. 根据 Fatou 引理,
\begin{align*}
\int_{X}\liminf_{n\to\infty}h_{n}\,d\mu\leqslant\liminf_{n\to\infty}\int_{X}h_{n}\,d\mu.
\end{align*}
但 $\liminf_{n\to\infty}h_{n}=g-f$ 且 $\int_{X}h_{n}\,d\mu=\int_{X}g\,d\mu-\int_{X}f_{n}\,d\mu$. 因此
\begin{align*}
\int_{X}g\,d\mu-\int_{X}f\,d\mu\leqslant\int_{X}g\,d\mu-\limsup_{n\to\infty}\int_{X}f_{n}\,d\mu.
\end{align*}
消去 $\int_{X}g\,d\mu$ 得到 $\limsup_{n\to\infty}\int_{X}f_{n}\,d\mu\leqslant\int_{X}f\,d\mu$.

类似地, 对 $k_{n}=g+f_{n}$ 应用 Fatou 引理:
\begin{align*}
\int_{X}g\,d\mu+\int_{X}f\,d\mu\leqslant\liminf_{n\to\infty}\left[\int_{X}g\,d\mu+\int_{X}f_{n}\,d\mu\right]=\int_{X}g\,d\mu+\liminf_{n\to\infty}\int_{X}f_{n}\,d\mu.
\end{align*}
因此 $\int_{X}f\,d\mu\leqslant\liminf_{n\to\infty}\int_{X}f_{n}\,d\mu$. 将两个不等式联立得到\eqref{align:ch18-dct}.
\end{proof}

\begin{corollary}\label{corollary:ch18-dct-bounded}
令 $(X,\mathcal{M},\mu)$ 为有限测度空间, 而 $\{f_{n}\}$ 是 $X$ 上的一致有界可测函数序列, 满足在 $X$ 上 a.e. 逐点地 $\{f_{n}\}\to f$. 则
\begin{align*}
\lim_{n\to\infty}\int_{X}f_{n}\,d\mu=\int_{X}f\,d\mu.
\end{align*}
\end{corollary}

\begin{proof}
若对所有 $n$, $|f_{n}|\leqslant M$ a.e., 则常函数 $g=M$ 是可积的, 因为 $\mu(X)<\infty$. 控制收敛定理适用.
\end{proof}

我们现在证明 Vitali 收敛定理, 它对一般测度空间成立.

\begin{definition}[一致可积与紧性]\label{definition:ch18-uniform-integrability}
令 $(X,\mathcal{M},\mu)$ 为测度空间, 而 $\{f_{n}\}$ 是 $X$ 上的可积函数序列. 序列 $\{f_{n}\}$ 称为在 $X$ 上一致可积的, 若对每个 $\varepsilon>0$, 存在 $\delta>0$, 使得对任何自然数 $n$ 和 $X$ 的可测子集 $E$,
\begin{align*}
\text{若 }\mu(E)<\delta,\quad\text{则 }\int_{E}|f_{n}|\,d\mu<\varepsilon.
\end{align*}
序列 $\{f_{n}\}$ 称为在 $X$ 上是紧的, 若对每个 $\varepsilon>0$, 存在 $X$ 的有限测度子集 $X_{0}$, 使得对任何自然数 $n$,
\begin{align*}
\int_{X\sim X_{0}}|f_{n}|\,d\mu<\varepsilon.
\end{align*}
\end{definition}

\begin{proposition}\label{proposition:ch18-integrability-criterion}
令 $(X,\mathcal{M},\mu)$ 为测度空间, 而 $f$ 是 $X$ 上的可积函数. 则对每个 $\varepsilon>0$, 存在 $\delta>0$, 使得对 $X$ 的任何可测子集 $E$,
\begin{align}
\text{若 }\mu(E)<\delta,\quad\text{则 }\int_{E}|f|\,d\mu<\varepsilon.\label{align:ch18-abs-continuity}
\end{align}
此外, 对每个 $\varepsilon>0$, 存在 $X$ 的有限测度子集 $X_{0}$, 使得
\begin{align}
\int_{X\sim X_{0}}|f|\,d\mu<\varepsilon.\label{align:ch18-tightness}
\end{align}
\end{proposition}

\begin{proof}
假设在 $X$ 上 $f\geqslant 0$. 令 $\varepsilon>0$. 由于 $\int_{X}f\,d\mu<\infty$, 存在简单函数 $\psi$ 满足 $0\leqslant\psi\leqslant f$ 且 $\int_{X}f\,d\mu-\int_{X}\psi\,d\mu<\varepsilon/2$. 选取 $M>0$ 使得在 $X$ 上 $0\leqslant\psi\leqslant M$. 对可测集 $E$,
\begin{align*}
\int_{E}f\,d\mu=\int_{E}\psi\,d\mu+\int_{E}(f-\psi)\,d\mu\leqslant M\cdot\mu(E)+\varepsilon/2.
\end{align*}
取 $\delta=\varepsilon/(2M)$, 若 $\mu(E)<\delta$, 则 $\int_{E}f\,d\mu<\varepsilon$.

对于紧性, 可测集 $X_{0}=\{x\in X\mid\psi(x)>0\}$ 有有限测度. 此外,
\begin{align*}
\int_{X\sim X_{0}}f\,d\mu=\int_{X\sim X_{0}}(f-\psi)\,d\mu\leqslant\int_{X}(f-\psi)\,d\mu<\varepsilon.
\end{align*}
一般情形从正部与负部得到.
\end{proof}

\begin{theorem}[Vitali 收敛定理]\label{theorem:ch18-vitali}
令 $(X,\mathcal{M},\mu)$ 为测度空间, 而 $\{f_{n}\}$ 是 $X$ 上的函数序列, 它在 $X$ 上既是一致可积的又是紧的. 假设在 $X$ 上 a.e. 逐点地 $\{f_{n}\}\to f$, 且 $f$ 在 $X$ 上可积. 则
\begin{align*}
\lim_{n\to\infty}\int_{X}f_{n}\,d\mu=\int_{X}f\,d\mu.
\end{align*}
\end{theorem}

\begin{proof}
对所有 $n$, 在 $X$ 上 a.e. $|f-f_{n}|\leqslant|f|+|f_{n}|$. 令 $\varepsilon>0$. 根据前一个命题, $\{f\}$ 是一致可积且紧的, 因此 $\{f_{n}-f\}$ 也是一致可积且紧的. 选取有限测度集 $X_{0}$ 使得对所有 $n$, $\int_{X\sim X_{0}}|f_{n}-f|\,d\mu<\varepsilon/3$. 在有限测度集 $X_{0}$ 上, 根据 Egoroff 定理, 存在 $X_{0}$ 的可测子集 $X_{1}$, 使得 $\mu(X_{0}\sim X_{1})<\delta$ 且在 $X_{1}$ 上 $\{f_{n}\}$ 一致收敛于 $f$. 选取 $\delta$ 使得对所有 $n$, $\int_{X_{0}\sim X_{1}}|f_{n}-f|\,d\mu<\varepsilon/3$. 对足够大的 $n$, 在 $X_{1}$ 上 $|f_{n}-f|<\varepsilon/[3\mu(X_{0})]$. 因此
\begin{align*}
\left|\int_{X}(f_{n}-f)\,d\mu\right|&\leqslant\int_{X_{1}}|f_{n}-f|\,d\mu+\int_{X_{0}\sim X_{1}}|f_{n}-f|\,d\mu+\int_{X\sim X_{0}}|f_{n}-f|\,d\mu\\
&<\varepsilon/3+\varepsilon/3+\varepsilon/3=\varepsilon.
\end{align*}
\end{proof}

\begin{corollary}\label{corollary:ch18-vitali-to-dct}
控制收敛定理是 Vitali 收敛定理的推论.
\end{corollary}

\begin{proof}
若 $|f_{n}|\leqslant g$ 其中 $g$ 可积, 则 $\{f_{n}\}$ 是一致可积且紧的, 因为 $\{g\}$ 是一致可积且紧的(命题\cref{proposition:ch18-integrability-criterion}).
\end{proof}

\begin{exercise}\label{exercise:ch18-23}
令 $(X,\mathcal{M},\mu)$ 为测度空间, 而 $f$ 是 $X$ 上的可积函数. 证明对每个 $\varepsilon>0$, 存在简单函数 $\psi$, 使得 $\int_{X}|f-\psi|\,d\mu<\varepsilon$.
\end{exercise}

\begin{exercise}\label{exercise:ch18-24}
证明: 若 $f$ 在 $X$ 上可积且 $E$ 是可测集, 则 $\int_{E}f\,d\mu=0$ 当且仅当在 $E$ 上 a.e. $f=0$.
\end{exercise}

\begin{exercise}\label{exercise:ch18-25}
令 $\{f_{n}\}$ 为 $X$ 上的非负可积函数序列, 满足在 $X$ 上 a.e. $f_{n}\to 0$. 证明 $\lim_{n\to\infty}\int_{X}f_{n}\,d\mu=0$ 当且仅当 $\{f_{n}\}$ 是一致可积且紧的.
\end{exercise}

\begin{exercise}\label{exercise:ch18-26}
令 $(X,\mathcal{M},\mu)$ 为测度空间, 而 $\{f_{n}\}$ 是 $X$ 上的可积函数序列. 证明: $\{f_{n}\}$ 是一致可积的当且仅当对每个 $\varepsilon>0$, 存在 $\delta>0$, 使得对任何自然数 $n$ 和可测集 $E$, 若 $\mu(E)<\delta$, 则 $\left|\int_{E}f_{n}\,d\mu\right|<\varepsilon$.
\end{exercise}

\begin{exercise}\label{exercise:ch18-27}
令 $\nu$ 为 $(X,\mathcal{M})$ 上的带号测度. 若 $f$ 在 $X$ 上关于 $\nu^{+}$ 和 $\nu^{-}$ 都可积, 定义 $f$ 关于 $\nu$ 的积分为
\begin{align*}
\int_{X}f\,d\nu=\int_{X}f\,d\nu^{+}-\int_{X}f\,d\nu^{-}.
\end{align*}
证明: 若在 $X$ 上 $|f|\leqslant M$, 则 $\left|\int_{X}f\,d\nu\right|\leqslant M|\nu|(X)$.
\end{exercise}

\begin{exercise}\label{exercise:ch18-28}
令 $g$ 为 $X$ 上的可积非负函数. 对所有 $E\in\mathcal{M}$, 定义 $\nu(E)=\int_{E}g\,d\mu$. 证明 $\nu$ 是 $(X,\mathcal{M})$ 上的测度. 此外, 对任何非负可测函数 $f$,
\begin{align*}
\int_{X}f\,d\nu=\int_{X}fg\,d\mu.
\end{align*}
（提示: 先对简单函数证明, 再用单调收敛定理.）
\end{exercise}

\end{document}